\chapter{Guix in the Wild: Domain Impact and Taming Real Hardware with Nonguix}
\label{chap:domain-impact-and-nonguix}

\epigraph{``In theory, there is no difference between theory and practice. In practice, your laptop's Wi-Fi card requires a proprietary binary blob, your university cluster has no root access, and your computational biology pipeline broke because an R package updated at midnight.''}{--- The Pragmatic Guix Alchemist}

\noindent
So far, we have explored GNU Guix from the perspective of computer science elegance: immutable stores, mathematical functions, declarative operating systems, and Scheme macros. But how does this functional paradigm fare when thrown into the messy, chaotic, real world?

In this chapter, we explore two transformative aspects of Guix in the wild:
\begin{enumerate}
  \item \textbf{Domain Impact}: How Guix is revolutionizing High-Performance Computing (HPC), bioinformatics, artificial intelligence, and academic research.
  \item \textbf{Pragmatic Hardware with Nonguix}: How to tame consumer laptops, proprietary Wi-Fi chips, NVIDIA GPUs, and non-free software without sacrificing a single drop of declarative purity.
\end{enumerate}

\section{Guix in Science \& Industry: The Reproducibility Revolution}

\subsection{1. High-Performance Computing (HPC) \& Supercomputing}

If you have ever used an institutional supercomputer or shared HPC cluster, you know the traditional pain:
\begin{itemize}
  \item You do not have \texttt{sudo} permissions (for good reason).
  \item The cluster runs an ancient distribution (like Red Hat Enterprise Linux 7 or CentOS) with GCC versions from a decade ago.
  \item You need bleeding-edge MPI libraries, specific CUDA toolkits, and dozens of custom packages.
  \item Traditional solutions like Environment Modules or Conda pollute paths, fail to isolate dynamic libraries, and break when compute nodes have slight architectural drifts.
\end{itemize}

\begin{alchemybox}[Unprivileged Cluster Freedom with Guix HPC]
Because GNU Guix package management is entirely unprivileged:
\begin{enumerate}
  \item Cluster sysadmins install the \texttt{guix-daemon} once.
  \item Hundreds of researchers can independently install, upgrade, and mix conflicting software stacks in their own profiles without ever filing an IT helpdesk ticket!
  \item Guix integrates natively with \textbf{SLURM}, allowing researchers to run batch jobs with \guixcmd{guix shell --container --manifest=manifest.scm -- srun ./simulation}.
\end{enumerate}
\end{alchemybox}

Leading research institutes---such as Inria in France, the Max Delbr\"uck Center in Germany, and the University of Tennessee---now manage multi-petaflop clusters powered by Guix HPC.

\subsection{2. Bioinformatics \& Computational Biology}

Bioinformatics pipelines are notorious for being the most fragile software ecosystems on the planet. A typical genomic pipeline might combine:
\begin{itemize}
  \item 40 different command-line tools written in C, C++, and Fortran.
  \item 25 Python packages with pinned C extensions.
  \item 150 R packages from CRAN and Bioconductor.
  \item A workflow engine like Nextflow or Snakemake.
\end{itemize}

A single minor version bump in an underlying alignment library can silently change variant calling results, leading to retracted scientific papers and millions of wasted research dollars.

With GNU Guix, bioinformaticians define their complete software environment in a single \texttt{manifest.scm} and pin the exact channel commits with \texttt{channels.scm}. Ten years later, any laboratory on Earth can re-run that pipeline with bit-for-bit identical results.

\begin{tearsbox}[The Replication Crisis in Computational Science]
In 2019, a major study attempted to reproduce the computational analysis of over 400 published machine learning and bioinformatics papers. Less than 25\% could be re-executed because of missing dependencies, changed numerical libraries, and unstated build flags.

Guix transforms reproducibility from an aspirational dream into an automated mathematical guarantee.
\end{tearsbox}

\subsection{3. Data Science, Machine Learning \& AI}

Modern AI engineering often feels less like programming and more like wrestling with dynamic linker errors:
\begin{center}
\texttt{ImportError: libcudart.so.11.0: cannot open shared object file}
\end{center}

Between CUDA versions, CuDNN drivers, Python wheels compiled against conflicting GLIBC versions, and PyTorch/TensorFlow bindings, ML developers frequently resort to 15GB monolithic Docker containers that nobody knows how to rebuild from source.

The \textbf{Guix Science} channel provides cleanly packaged, transparently built PyTorch, JAX, and CUDA/ROCm stacks where every GPU acceleration library is explicitly tracked in the dependency graph.

\subsection{4. Verifiable Software Supply Chains (SLSA Level 4)}

In enterprise software engineering, security attacks on the build pipeline (such as the SolarWinds compromise) have made verifiable builds a critical priority.

Because Guix builds are \textbf{hermetic}, timestamp-independent, and isolated from host impurities, Guix achieves industry-leading \textbf{Bootstrappable and Reproducible Builds}:
\begin{itemize}
  \item You can independently compile a package on three different machines across the world and verify that they produce the exact, bit-for-bit identical SHA-256 binary hash.
  \item Guix can bootstrap its entire C compiler toolchain starting from a tiny 357-byte binary seed (\texttt{stage0}), completely eliminating compiler backdoors (Ken Thompson's ``Trusting Trust'' attack)!
\end{itemize}

\section{The Imperfect World: Consumer Hardware \& The Driver Dilemma}

GNU Guix is maintained by the GNU Project, which means its default configuration adheres strictly to the \textbf{GNU Free System Distribution Guidelines (FSDG)}:
\begin{itemize}
  \item The default kernel is \textbf{Linux-Libre} (the Linux kernel stripped of all non-free binary firmware blobs).
  \item No proprietary graphics drivers (like NVIDIA's proprietary driver) or proprietary commercial apps are included in official channels.
\end{itemize}

\subsection{The Reality of Modern Laptops}

While this ethical purity is noble, modern consumer hardware is often hostile to free software:
\begin{itemize}
  \item \textbf{Wi-Fi \& Bluetooth}: Almost all modern Intel (AX200/AX211), Realtek (RTL8821), and Broadcom wireless chips require proprietary firmware blobs loaded onto the chip at boot time. Without these blobs, the operating system cannot even see the Wi-Fi card.
  \item \textbf{Graphics}: NVIDIA GeForce and RTX discrete GPUs require proprietary firmware and kernel modules for hardware 3D acceleration, CUDA compute, and external display output.
  \item \textbf{CPU Microcode}: Modern Intel and AMD x86 processors contain microcode vulnerabilities (Spectre, Meltdown, Retbleed, Downfall, Zenbleed). Mitigating these requires loading proprietary CPU microcode updates at early boot.
  \item \textbf{Workplace Software}: Corporate communication tools like Slack, Zoom, Microsoft Teams, and Google Chrome are closed-source.
\end{itemize}

If you attempt to install pure Guix System on a modern gaming laptop or ultrabook without proprietary firmware, you may find yourself with no Wi-Fi, no GPU acceleration, and an unaccelerated software display renderer.

\section{Nonguix: The Pragmatic Bridge}

To solve this dilemma without compromising the core project's principles, the community created \textbf{Nonguix} (\url{https://gitlab.com/nonguix/nonguix}).

Nonguix is an independent, decentralized channel that packages:
\begin{enumerate}
  \item \textbf{Vanilla Linux Kernel (\texttt{linux})}: The full upstream Linux kernel containing all driver blobs.
  \item \textbf{\texttt{linux-firmware}}: Binary firmware packages for Intel, AMD, Realtek, MediaTek, and Broadcom.
  \item \textbf{\texttt{microcode-initrd}}: Early-boot CPU microcode loader for Intel and AMD processors.
  \item \textbf{Proprietary Software}: Steam, Spotify, Google Chrome, Discord, Slack, and NVIDIA drivers.
  \item \textbf{Pre-Built Binary Substitutes}: A dedicated build farm at \texttt{https://substitutes.nonguix.org} so you never have to wait hours compiling the Linux kernel or Chromium on your laptop!
\end{enumerate}

\begin{footgunbox}[Compiling the Linux Kernel by Mistake]
If you use Nonguix without authorizing its pre-built substitute server, Guix will happily download the 150MB Linux kernel source tarball and compile it on your machine, spinning your laptop fans at 100\% for 45 minutes! Always authorize \texttt{substitutes.nonguix.org} before running system reconfiguration.
\end{footgunbox}

\section{Case Study: The Ultimate Real-World Workstation / Laptop}

Let us bring everything together with a complete, production-ready \texttt{/etc/config.scm} for a modern laptop or desktop with Wi-Fi, CPU microcode updates, sound (PipeWire), and power management:

\begin{lstlisting}[style=guixscheme]
;; /etc/config.scm - Battle-tested real-world laptop configuration
(use-modules (gnu)
             (gnu system nss)
             ;; Nonguix modules for kernel and firmware
             (nongnu packages linux)
             (nongnu system linux-initrd))

(use-service-modules desktop networking ssh cups sound)
(use-package-modules certs fonts gnome linux package-management version-control)

(operating-system
  (host-name "nomad-thinkpad")
  (timezone "America/New_York")
  (locale "en_US.utf8")

  ;; 1. Use the non-libre Linux kernel with early CPU microcode updates
  (kernel linux)
  (initrd microcode-initrd)
  (firmware (list linux-firmware))

  ;; 2. Bootloader Configuration (UEFI GRUB)
  (bootloader (bootloader-configuration
                (bootloader grub-efi-bootloader)
                (targets '("/boot/efi"))
                (keyboard-layout keyboard-layout)))

  ;; 3. Storage and Filesystems
  (file-systems (cons* (file-system
                         (mount-point "/")
                         (device (file-system-label "guix-root"))
                         (type "ext4"))
                       (file-system
                         (mount-point "/boot/efi")
                         (device (uuid "ABCD-1234" 'fat))
                         (type "vfat"))
                       %base-file-systems))

  ;; 4. User Configuration
  (users (cons (user-account
                 (name "sudheer")
                 (group "users")
                 (supplementary-groups '("wheel" "netdev" "audio" "video" "kvm" "lp"))
                 (home-directory "/home/sudheer"))
               %base-user-accounts))

  ;; 5. System-wide Packages
  (packages (append (list nss-certs
                          font-dejavu
                          font-google-noto
                          git
                          emacs)
                    %base-packages))

  ;; 6. System Services & Nonguix Substitute Server
  (services
    (append
      (list (service gnome-desktop-service-type)
            (service bluetooth-service-type
              (bluetooth-configuration
                (auto-enable? #t))))

      ;; Modify default services to add Nonguix pre-built binary caches
      (modify-services %desktop-services
        (guix-service-type config =>
          (guix-configuration
            (inherit config)
            (substitute-urls
              (append (list "https://substitutes.nonguix.org"
                            "https://bordeaux.guix.gnu.org")
                      %default-substitute-urls))
            (authorized-keys
              (append (list (plain-file "nonguix-signing-key"
                              "(public-key
 (ecc
  (curve Ed25519)
  (q #C1FD53E5D4A3D2A0B388D378307D5E06294D2AC7F453EB278F36990C28A13D02#)))"))
                      %default-authorized-guix-keys))))))))
\end{lstlisting}

\section{The Best of Both Worlds}

With this architecture, you achieve what was once thought impossible in the Linux world:
\begin{enumerate}
  \item \textbf{100\% Real-World Hardware Compatibility}: Flawless Wi-Fi, Bluetooth audio, sleep/suspend, and GPU acceleration.
  \item \textbf{100\% Declarative Purity}: Your entire operating system is defined in one version-controlled Scheme file.
  \item \textbf{Zero Fear of Breakage}: Every update generates a new bootloader entry. If an experimental kernel update misbehaves, selecting Generation $N-1$ in GRUB returns your machine to working order in 10 seconds.
\end{enumerate}

Whether deploying computational clusters across hundreds of supercomputer nodes or running a sleek laptop workstation, GNU Guix provides the ultimate foundation for reliable computing.
